% === EXECUTIVE SUMMARY ===
\chapter{Executive Summary}

Group conversations often fail to achieve their potential. Participants interrupt each other, some voices dominate while others remain unheard, and discussions frequently drift off track~\cite{Diehl1987ProductionBlocking}. These dynamics undermine collaboration, making it difficult for groups to share knowledge effectively or reach clear decisions.

AIutami addresses these challenges through an AI-moderated discussion platform. The system allows users to create and join audio-based sessions through a web interface. Each session follows a structured workflow: participants gather in a lobby phase before entering the active discussion. During the conversation, speaking turns are regulated through a request-based mechanism that prevents overlap and ensures balanced participation. Visual feedback keeps participants informed about the current speaker and turn status, improving coordination and reducing confusion.

What distinguishes AIutami from traditional conferencing tools is its intelligent moderation layer. The platform combines real-time audio streaming, automatic speech recognition, and large language models to monitor conversational dynamics. When the discussion becomes unbalanced, stuck, or loses focus, the AI moderator intervenes contextually---prompting quieter participants, redirecting off-topic exchanges, or helping the group move forward.

Consider a Murder Mystery activity where participants must share clues and collaborate to identify the culprit. Without structure, a few voices dominate, clues get lost, and the group struggles to reason together. With AIutami, every participant gets heard, key points remain visible, and the group can build toward a shared conclusion. The same benefits apply to study groups debating complex topics, teams making strategic decisions, or workshops requiring structured input from all participants.

AIutami delivers concrete value: discussions that are fairer, more focused, and more likely to produce clear outcomes.
